\section{Conclusions}

A non-detection is an upper bound and an open-world statement. Recording it as a
zero replaces that statement with a number; routine data models offer nowhere
else to put it.

Across \num{4190833} river station-years, \num{44.9}\% of samples lie below the
quantification limit, and \num{25.5}\% of station-years carry neither a censoring
flag nor a limit, so a zero and a measured trace cannot be told apart. Where a
standard exists, \num{40.1}\% of assessments come from a method failing the
performance criterion of Directive 2009/90/EC, and \num{17.5}\% from one whose
limit exceeds the standard outright --- results Article~3(3b) requires to be set
aside and a two-valued pipeline records as compliance.

The record is not one era. The share declaring neither a flag nor a limit falls
from \num{49.7}\% before \num{2015} to \num{0.0}\% after, and the \num{8}
authorities reporting in both periods account for all of it. The
representational correction is work these reporters have already done, under an
unchanged standard. The analytical failure does not move: on \num{2024} alone,
with every limit reported, \num{47.7}\,\% of assessments still come from a
method that cannot decide its own standard.

Over \num{696168} assessments --- station-years, not samples --- the same
records yield a \num{3.4}-fold range in exceedances according to whether a
non-detection is entered at zero, half the limit or the limit; \num{14505}
are exceedances the law can affirm: the number a two-valued pipeline reports is
fixed by the convention, not by the chemistry. Annex~I also
states a second standard on a different unit of observation, and reading the
\num{96597294} individual samples of the disaggregated release leaves
\num{9.7}\,\% of them undecidable for the same reason, with \num{25.6}\,\% of
the station-years assessable under both standards receiving different verdicts
--- two answers a single threshold column cannot hold at once.

A two-valued model cannot express these states. CENSO can: it binds the
quantification limit to the result it bounds, and makes \emph{cannot be
determined} the third of the three outcomes the Directive admits --- subtyped by
the reason, so that an unmeasurable standard, an unreported bound, an
inapplicable one and an interval straddling the limit are four answers and not
one shrug. Its axioms reject all eight defects drawn
from this dataset and accept both controls, and expressing the record in it
loses nothing: every row classified in the graph lands where the streaming
count puts it. Applying two jurisdictions to
the same observations changes \num{17.7}\% of the outcomes both regulate, with
no code change. Of the \num{21} ontologies parsed, none combines these
capabilities and none carries a censoring status at all, and the shortfall is
not lexical: on one real record \num{20} of them cannot separate a non-detection
from a measurement reporting the same number,
because the two readings translate to the same facts in each. CENSO makes them
incompatible rather than merely differently labelled.

The corrective for practice is cheap: record the quantification limit with every
result, per analytical method. The corrective for representation is to stop
treating ``not detected'' as a number.